%%%%%%%%%%%%%%%%%%%%%%%%%%%%%%%%%%%%%%%%%%%%%%%%%%%%%%%%%%%%%%%%% 3. METRICS
\section{Evaluation Metrics and Their Interpretation}

Six metrics are computed on every model. This section explains each one, how to
read its value, and --- critically --- whether any standard threshold actually
exists.

\begin{warnbox}[Read this before quoting any threshold]
There is \textbf{no universally accepted ``good'' value} for DPD, EOD, or GE.
This guide distinguishes three levels of authority:
\begin{itemize}
\item \textbf{Established standard} --- a formal, citable rule. Almost nothing
in fairness metrics qualifies.
\item \textbf{Common rule of thumb} --- widely used in practice but with no
formal backing. Label it as such in the paper.
\item \textbf{Project-specific interpretation} --- our own reading of our own
numbers.
\end{itemize}
\end{warnbox}

\subsection{Accuracy}

\textbf{What:} fraction of correct predictions.
\textbf{Direction:} higher is better.
\textbf{Range:} $0$ to $1$.

\textbf{Reference point:} the majority-class rate. Base rates in this project
are COMPAS $0.455$ (majority-class $\approx 0.545$), German $0.300$
($\approx 0.70$), Adult $0.248$ ($\approx 0.752$). A model must beat that to be
doing anything at all.

\textbf{Standard?} Accuracy itself is universal; there is \textbf{no universal
``good'' accuracy}. What exists is \emph{literature-typical} performance per
dataset --- COMPAS models land at $\approx 0.66$--$0.70$, Adult
$\approx 0.84$--$0.87$, German $\approx 0.72$--$0.77$.

\subsection{DPD --- Demographic Parity Difference}

\textbf{Stands for:} Demographic Parity Difference (also: statistical parity
difference).

\textbf{Measures:} the largest gap in \textbf{selection rate} (fraction
predicted positive) between sensitive groups:
\begin{equation}
\mathrm{DPD} = \max_g P(\hat y = 1 \mid S = g) - \min_g P(\hat y = 1 \mid S = g)
\end{equation}
It \textbf{ignores the true labels entirely} --- it only asks who gets flagged,
not who deserved it.

\textbf{Direction:} lower is better. \textbf{Range:} $0$ to $1$. Zero means
every group is flagged at an identical rate.

\begin{warnbox}[Is there a standard threshold? No.]
What people \emph{do} cite:
\begin{itemize}
\item \textbf{The four-fifths rule / 80\% rule} (US EEOC, \emph{Uniform
Guidelines on Employee Selection Procedures}, 1978) --- a legal rule of thumb
for adverse impact. But it is a \textbf{ratio} test
($\min\text{ rate}/\max\text{ rate} \ge 0.8$), \textbf{not a difference test}.
It does not translate to a DPD threshold without knowing the base selection
rate. Cite it for context; do \emph{not} claim ``DPD $<0.2$ passes the 80\%
rule'' --- that is only true if the maximum selection rate happens to be $1.0$.
\item \textbf{Common informal practice:} DPD $<0.05$ is treated as ``close to
parity,'' $0.05$--$0.10$ as modest disparity, $>0.10$ as substantial. This is a
\emph{convention}, not a standard.
\end{itemize}
\end{warnbox}

\subsubsection{The two variants in this project}
\begin{itemize}
\item \code{Demographic Parity Difference} --- computed over \textbf{all} groups.
\item \code{DPD (Largest 2 Groups)}, written \dpdtwo{} --- restricted to the two
largest groups. On COMPAS that is African-American ($n=3175$) vs Caucasian
($n=2103$), ProPublica's focal comparison. On German and Adult (two groups
each) the two variants are \textbf{identical by construction}.
\end{itemize}

\begin{example}[Why the largest-2 variant exists]
COMPAS has six race groups and the small ones are tiny: Asian $n=31$, Native
American $n=11$ across the \emph{whole} dataset --- roughly $6$ and $2$ people
in a 1{,}235-row test split. A $\max-\min$ over all groups is then pinned by a
handful of individuals and swings wildly. The evidence is in our own results:
\begin{center}
\begin{tabular}{@{}lcc@{}}
\toprule
COMPAS, LR baseline & mean & std \\
\midrule
DPD (all six groups)  & $0.608$ & $\pm 0.211$ \\
\dpdtwo{} (two largest) & $0.304$ & $\pm 0.049$ \\
\bottomrule
\end{tabular}
\end{center}
The standard deviation is over four times larger for the all-groups version.
That is why \dpdtwo{} is the primary metric --- and this table \emph{is} the
justification to put in the paper.
\end{example}

\subsection{EOD --- Equalized Odds Difference}

\textbf{Measures:} whether the model's \emph{error rates} are equal across
groups. Computed as the larger of \{max$-$min true-positive-rate gap,
max$-$min false-positive-rate gap\}. Unlike DPD it \emph{does} use the true
labels --- it asks whether the model is equally accurate for each group, not
whether it selects them equally often.

\textbf{Direction:} lower is better. \textbf{Range:} $0$ to $1$.
\textbf{Standard threshold?} \textbf{No} --- same situation as DPD.

\begin{keybox}[The relationship that drives our biggest result]
DPD and EOD are \textbf{mathematically incompatible} whenever base rates differ
across groups. Our data has exactly that: Adult female base rate $0.114$ vs
male $0.312$; COMPAS African-American $0.523$ vs Caucasian $0.391$. So
\textbf{driving DPD to zero must push EOD up.} Our Adult results show this
precisely.
\end{keybox}

\begin{warnbox}[COMPAS EOD is unreliable in the current code]
EOD is computed over \textbf{all} groups, without the largest-2 restriction
that was applied to DPD. So COMPAS EOD inherits exactly the tiny-group noise
problem --- which is why every COMPAS EOD is $\approx 0.64$--$0.83$ with std
$\approx 0.15$--$0.22$, and why \emph{no} COMPAS EOD comparison reaches
significance. Either add a largest-2 EOD, or explicitly caveat the column. Do
not report it as a clean finding.
\end{warnbox}

\subsection{GE(2) --- Generalized Entropy Index (labelled ``Theil'' in the code)}

\textbf{What it is:} inequality of the benefit vector $b = 1 + \hat y - y$,
measured with the generalized entropy index of order $2$:
\begin{equation}
\getwo = \frac{1}{2n}\sum_i\left[\left(\frac{b_i}{\bar b}\right)^{2} - 1\right]
       = \frac{\mathrm{Var}(b)}{2\,\bar b^{\,2}}
\end{equation}

\textbf{Direction:} lower is better. \textbf{Range:} $0$ to $\infty$ in
principle; $0$ means perfectly equal benefit (a perfect classifier, since every
$b=1$).

\begin{warnbox}[Naming error to fix before submitting]
\file{IndividualFairness.py} defines \code{theil\_index} as AIF360's
\code{generalized\_entropy\_error(..., alpha=2)}. That is \textbf{backwards
relative to the standard definition}. AIF360's own docstring states: a value of
$0$ is the mean log deviation, \textbf{$1$ is the Theil index}, and $2$ is half
the squared coefficient of variation.

So everything the project calls ``Theil Index'' is really \textbf{GE($\alpha=2$)}.
\textbf{No number is wrong} --- GE(2) is a perfectly standard inequality index,
and crucially it is the \emph{same} index the SoftGE surrogate optimises (also
$\alpha=2$), so training and evaluation are consistent. But the \textbf{label}
must change, or a reviewer familiar with Speicher et al.\ will flag it
immediately. The paper produced alongside this guide already uses
\textbf{GE(2)} throughout.
\end{warnbox}

\textbf{What it actually captures on hard predictions:} with binary
predictions, $b \in \{0,1,2\}$ only. So GE(2) here is a \textbf{function of the
false-positive and false-negative counts alone}. It measures how much error
there is and how it splits between the two error types --- not, strictly, whether
\emph{similar individuals} got similar treatment. This is the standard Speicher
et al.\ operationalisation, so it is defensible, but be precise: it is an
\emph{inequality-of-outcomes} measure, not a Lipschitz-style ``similar
individuals treated similarly'' measure.

\textbf{Standard threshold?} \textbf{None whatsoever.} GE has no absolute
interpretation. It is meaningful \textbf{only comparatively} --- model A versus
model B on the same dataset. Say this plainly in the paper rather than implying
a range.

\subsection{Gini Coefficient}

\textbf{What it is supposed to be:} the classic economics inequality measure,
from the Lorenz curve. $0 =$ perfect equality, $1 =$ maximum inequality.

\begin{warnbox}[This metric carries no fairness information]
\file{IndividualFairness.py} applies the Lorenz/Gini construction to
\code{y\_pred} directly. Since \code{y\_pred} is a vector of 0s and 1s, this
reduces \textbf{exactly} to
\[ \mathrm{Gini} = 1 - \text{positive prediction rate} \]
Verified numerically to four decimal places on three random vectors (positive
rate $0.3270 \to$ Gini $0.6730$; $0.2045 \to 0.7955$; $0.4700 \to 0.5300$), and
consistent with the result tables (German LR baseline Gini $0.869$ implies a
positive rate of $0.131$ against a base rate of $0.30$ --- plausible
under-prediction on an imbalanced dataset).

\textbf{Recommendation:} either \textbf{drop it} from the paper, or
\textbf{relabel it} as ``selection rate (reported as $1-$Gini)'' --- which is
genuinely useful, since it is exactly the quantity the degeneracy guard
watches. What you must not do is present it as an independent
individual-fairness metric.
\end{warnbox}

\subsection{Soft DP / Soft GE}

Logged per sweep point, computed on \emph{test probabilities}. These are not
evaluation metrics --- they exist solely for the surrogate-validation scatter
(Figure~\ref{fig:guide-surrogate}). Keep them in supplementary material.

\subsection{Summary table}

\begin{center}
\scriptsize
\renewcommand{\arraystretch}{1.3}
\setlength{\tabcolsep}{3.5pt}
\begin{tabular}{@{}L{0.098\textwidth} L{0.205\textwidth} c L{0.055\textwidth} L{0.195\textwidth} L{0.245\textwidth}@{}}
\toprule
\textbf{Metric} & \textbf{Meaning (plain)} & \textbf{Dir.} & \textbf{Range}
& \textbf{Standard?} & \textbf{How we interpret our result} \\
\midrule

\textbf{Accuracy} & \% correct & $\uparrow$ & $0$--$1$ &
No. Compare to majority-class rate and dataset literature &
COMPAS ${\approx}0.68$, Adult ${\approx}0.85$, German ${\approx}0.73$ --- all
literature-typical \\

\textbf{DPD} (all groups) & Max$-$min selection rate over all groups &
$\downarrow$ & $0$--$1$ &
\textbf{No.} 80\% rule is a \emph{ratio} test, not a difference test &
On COMPAS dominated by tiny-group noise --- report as secondary only \\

\textbf{\dpdtwo{}} & Same, restricted to the two biggest groups &
$\downarrow$ & $0$--$1$ &
\textbf{No.} ``${<}0.05$ = near parity'' is convention, not standard &
\textbf{Primary group-fairness metric.} Stable across seeds \\

\textbf{EOD} & Max gap in TPR \emph{or} FPR across groups &
$\downarrow$ & $0$--$1$ &
\textbf{No.} ``${<}0.1$'' is convention only &
Mathematically opposed to DPD when base rates differ.
\textbf{COMPAS values unreliable} \\

\textbf{\getwo{}} & Inequality of benefit $b = 1{+}\hat y{-}y$ &
$\downarrow$ & $0$--$\infty$ &
\textbf{None. Comparative only} &
\textbf{Primary individual-fairness metric.} Only meaningful vs another model
on the same data \\

\textbf{Gini} & \emph{As implemented:} exactly $1-$selection rate &
--- & $0$--$1$ & n/a &
\textbf{Not an independent fairness metric.} Drop or relabel \\

\textbf{Soft DP / Soft GE} & Differentiable training surrogates &
$\downarrow$ & --- & n/a &
Supplementary: validates soft ${\approx}$ hard ($r = 0.63$--$0.92$) \\

\bottomrule
\end{tabular}
\\[4pt]
\footnotesize $\uparrow$ = higher is better, $\downarrow$ = lower is better.
\end{center}

%%%%%%%%%%%%%%%%%%%%%%%%%%%%%%%%%%%%%%%%%%%%%%%%%%%%%%%%%%%%% 4. METHODOLOGY
\section{Methodology, Step by Step}

\subsection{Step 1 --- Dataset selection and loading}

Three datasets, chosen because together they are the field-standard trio for
tabular fairness work. Actual figures, read from the loaders:

\begin{center}
\small
\begin{tabular}{@{}lccc@{}}
\toprule
 & \textbf{COMPAS} & \textbf{German Credit} & \textbf{Adult} \\
\midrule
Samples $n$            & $6{,}172$ & $1{,}000$ & $45{,}222$ \\
Target ($y=1$)         & re-offended & \textbf{bad} credit & income $>$\$50K \\
Sensitive attribute    & race & Sex & sex \\
Number of groups       & $6$ & $2$ & $2$ \\
Overall base rate      & $0.455$ & $0.300$ & $0.248$ \\
Encoded features       & $12$ & $21$ & $\approx 79$ \\
\midrule
\multicolumn{4}{@{}l}{\textit{Group sizes and per-group base rates}} \\
Largest group  & Afr.-Am.\ $3{,}175$ ($0.523$) & male $690$ ($0.277$) & Male $30{,}527$ ($0.312$) \\
2nd largest    & Caucasian $2{,}103$ ($0.391$) & female $310$ ($0.352$) & Female $14{,}695$ ($0.114$) \\
Others         & Hisp.\ $509$, Other $343$, & --- & --- \\
               & \textbf{Asian 31, Nat.Am.\ 11} & & \\
\bottomrule
\end{tabular}
\end{center}

\textbf{Why these three:} they span three domains (criminal justice, credit,
income), two sensitive attributes (race, sex), and three orders of magnitude of
sample size --- exactly what RQ2 needs to test consistency.

\textbf{Why the differing base rates matter:} they are the \emph{reason} DPD and
EOD conflict. Adult's female-vs-male base rate gap ($0.114$ vs $0.312$, a factor
of $2.7$) is the mechanical driver of our most striking result.

\subsection{Step 2 --- The COMPAS leakage fix}

We apply ProPublica's standard row filters:
\begin{itemize}
\item \code{days\_b\_screening\_arrest} within $\pm30$ days
\item \code{is\_recid} $\ne -1$
\item \code{c\_charge\_degree} $\ne$ \code{O}
\item \code{score\_text} $\ne$ \code{N/A}
\end{itemize}
Features used:
\code{age}, \code{sex}, \code{juv\_fel\_count}, \code{juv\_misd\_count},
\code{juv\_other\_count}, \code{priors\_count}, \code{c\_charge\_degree},
\code{race}.

\begin{keybox}[This is a genuine methodological contribution]
An earlier version of this project included a feature
$\code{duration} = \code{end} - \code{start}$ that correlated $\boldsymbol{-0.78}$
with the target. That correlation is \textbf{mechanical, not predictive}:
re-offending truncates the observation window, so short duration $\Rightarrow$
re-offended, almost by definition. That is \textbf{label leakage} --- the model
was reading the answer off the feature.

Removing it drops COMPAS accuracy to the literature-typical
$\approx 0.67$--$0.68$. \textbf{This is a correction, not a regression.} Frame it
as a soundness contribution in the paper.
\end{keybox}

\subsection{Step 3 --- Preprocessing}

Per seed:
\begin{enumerate}
\item Split $60/20/20$ into train/validation/test. (Note: splits are
\emph{not} stratified --- worth a sentence in the paper.)
\item Extract the sensitive column per split; build integer group codes with
\textbf{one shared category mapping across all splits} so group IDs mean the
same thing everywhere. This matters --- the torch loss indexes by these codes.
\item Fit the \code{ColumnTransformer} \textbf{on train only}, apply to
val/test --- correct, no leakage:
  \begin{itemize}
  \item numeric $\to$ \code{StandardScaler}
  \item categorical $\to$ \code{OneHotEncoder(drop='first', handle\_unknown='ignore')}
  \end{itemize}
\item Convert to torch tensors.
\end{enumerate}

\begin{warnbox}[There is no feature-engineering step --- say so]
No interaction terms, no binning, no feature selection. Features are used as
published. This is \emph{appropriate} --- the paper is about the loss function,
and inventing features would confound the comparison. But state it explicitly
rather than leaving a gap where a ``feature engineering'' section would be.
\end{warnbox}

\subsection{Step 4 --- The composite loss}

Equation~\eqref{eq:composite}, with SoftDP and SoftGE as defined earlier. Two
bugs from the legacy code are fixed here, and both belong in the paper's
reproducibility discussion:

\begin{enumerate}
\item \textbf{Gradient flow.} The original computed fairness terms via
fairlearn/aif360 on the detached array
\begin{center}\code{y\_pred.round().detach().numpy()}\end{center}
The \code{.detach()} severs the autograd graph --- \textbf{only BCE ever
produced gradients}, so $\alpha$ and $\beta$ had almost no effect on the learned model.
That is why the legacy heatmaps were flat in $\beta$. Every term is now computed
in torch on raw probabilities.
\item \textbf{Double sigmoid.} The old model applied sigmoid in \code{forward}
and the loss applied it again. Now models return logits and the loss uses
\code{binary\_cross\_entropy\_with\_logits}.
\end{enumerate}

\begin{example}[Putting a number through the composite loss]
Suppose at some point in training, on the training batch:
\[
\mathrm{BCE} = 0.62, \qquad \mathrm{SoftDP} = 0.18, \qquad \mathrm{SoftGE} = 0.09
\]
Compare two settings of the dial at the same fairness pressure
$\alpha = 0.75$ (so $1-\alpha = 0.25$):

\medskip
\textbf{$\beta = 1$ (group-only):}
\[
\mathcal{L} = 0.75(0.62) + 0.25\big[1(0.18) + 0(0.09)\big]
            = 0.465 + 0.045 = \mathbf{0.510}
\]

\textbf{$\beta = 0$ (individual-only):}
\[
\mathcal{L} = 0.75(0.62) + 0.25\big[0(0.18) + 1(0.09)\big]
            = 0.465 + 0.0225 = \mathbf{0.4875}
\]

\textbf{$\beta = 0.5$ (blend):}
\[
\mathcal{L} = 0.465 + 0.25\big[0.5(0.18) + 0.5(0.09)\big]
            = 0.465 + 0.03375 = \mathbf{0.4988}
\]

Same accuracy term in all three; only the \emph{composition} of the fairness
penalty changes. Gradient descent therefore pushes the model in a different
direction for each $\beta$ --- which is the entire mechanism the study
exploits.
\end{example}

\subsection{Step 5 --- Training}

Adam (lr $0.01$, weight decay $10^{-4}$), full-batch, up to 300 epochs, early
stopping on \textbf{validation composite loss} with patience 20, best weights
restored from a deep copy.

Note that early stopping monitors the \emph{composite} loss at that
$(\alpha,\beta)$ --- not plain BCE. This is the right choice (you select the best
model for the objective you are actually optimising), but it means the stopping
criterion shifts as you move across the grid. Mention it.

\subsection{Step 6 --- The sweep}

For each of $7\alpha \times 7\beta = 49$ combinations, per architecture (2), per
dataset (3), per seed (10):

\begin{itemize}
\item Set \code{torch.manual\_seed(seed)} and \code{np.random.seed(seed)}
before instantiating the model $\Rightarrow$ identical initialisation across
grid points within a seed, so differences are attributable to $(\alpha,\beta)$,
not initialisation luck.
\item Train (or load a cached checkpoint).
\item Predict on test, threshold at $0.5$, compute all six metrics.
\item Also record validation accuracy, validation \dpdtwo{}, validation positive
rate (for selection and the degeneracy guard), soft metrics, and training time.
\end{itemize}

\begin{center}
\textbf{Total: $49 \times 2 \times 3 \times 10 = 2{,}940$ trained models.}
\end{center}

\subsection{Step 7 --- Model selection}

Validation-only, in three steps:
\begin{enumerate}
\item Keep $\alpha < 1$ (must have \emph{some} fairness weight).
\item Keep validation positive rate $\in [0.05, 0.95]$ (degeneracy guard).
\item Minimise $\sqrt{(1-\text{val\_acc})^2 + \text{val\_DPD}_2^2}$.
\item Test metrics reported only, \textbf{never used for selection}.
\end{enumerate}

Which $(\alpha, \beta)$ actually got picked, across ten seeds --- this is itself
a result:

\begin{center}
\small
\begin{tabular}{@{}llcc@{}}
\toprule
\textbf{Dataset} & \textbf{Model} & \textbf{Distinct configs / 10 seeds} & \textbf{Mean $\beta$} \\
\midrule
COMPAS & Fair LR  & $5$  & $0.63$ \\
COMPAS & Fair MLP & $5$  & $0.60$ \\
German & Fair LR  & $8$  & $0.19$ \\
German & Fair MLP & $\mathbf{10}$ & $0.41$ \\
Adult  & Fair LR  & $4$  & $0.58$ \\
Adult  & Fair MLP & $3$  & $0.65$ \\
\bottomrule
\end{tabular}
\end{center}

German is unstable: the MLP picks a \emph{different} configuration in every
single seed. That instability is a finding, analysed in
Section~\ref{sec:guide-german}.

\subsection{Step 8 --- Baselines and Step 9 --- Aggregation}

All baselines are fit on the same seed's training split and evaluated on the
same test split. Then the ten per-seed CSVs are concatenated, and for each
(fair model, baseline) pair we run a \textbf{paired Wilcoxon signed-rank test}
across seeds per metric, \textbf{Holm-corrected} across the metric family
\{Accuracy, \dpdtwo{}, EOD, \getwo{}\}.

\begin{warnbox}[A statistical detail to get right]
With $n=10$ seeds, the smallest two-sided $p$-value a Wilcoxon test can produce
is $2^{-9} \approx 0.00195$. Report it as $p = 0.002$. Do \textbf{not} write
``$p < 0.001$'' --- that is not attainable at this sample size, and it means
every single seed moved in the same direction.
\end{warnbox}

\subsection{Libraries used}

\begin{center}
\small
\begin{tabular}{@{}ll@{}}
\toprule
\textbf{Library} & \textbf{Used for} \\
\midrule
PyTorch & models + the differentiable composite loss \\
scikit-learn & preprocessing, splits, LR/RF, accuracy \\
fairlearn & DPD, EOD, ExponentiatedGradient, ThresholdOptimizer \\
AIF360 & generalized entropy index \\
xgboost & XGBoost accuracy reference \\
scipy & Wilcoxon signed-rank test \\
pandas / numpy & data handling \\
matplotlib / seaborn & figures \\
\bottomrule
\end{tabular}
\end{center}
